\documentclass[a4paper,14pt]{extarticle}
\usepackage[utf8]{inputenc}
\usepackage[T2A]{fontenc}
\usepackage[english,russian]{babel}
\usepackage{geometry}
\geometry{left=25mm, right=15mm, top=20mm, bottom=20mm}

\begin{document}

\begin{titlepage}
\begin{center}
\vspace*{2cm}
{\Large \textbf{ОПИСАНИЕ ПРОГРАММЫ}} \\
\vspace{0.5cm}
{\large \textbf{«Motion Visualizer»}} \\
\vspace{0.5cm}
\vfill
г. Ростов-на-Дону, 2026
\end{center}
\end{titlepage}

\section{Назначение программы}
Программа демонстрирует движение 3D-объекта (куба) в трёхмерном пространстве с регистрацией параметров в базе данных. Предназначена для демонстрации связки: GUI → C++ ядро → 3D-визуализация → БД.

\section{Функциональные характеристики}

\subsection{Основные функции}
\begin{enumerate}
    \item Запуск движения (кнопка «Старт»)
    \item Остановка и сброс (кнопка «Сброс»)
    \item Перетаскивание объекта мышью
    \item Просмотр истории испытаний (кнопка «Показать логи»)
    \item Автоматическое сохранение в SQLite
\end{enumerate}

\subsection{Входные данные}
\begin{itemize}
    \item Скорость движения (задана в коде: 50 px/сек)
    \item Положение курсора мыши (при перетаскивании)
\end{itemize}

\subsection{Выходные данные}
\begin{itemize}
    \item Координаты объекта на экране
    \item Записи в БД: ID, сценарий, время старта/финиша, дистанция, статус
\end{itemize}

\section{Условия эксплуатации}

\subsection{Требования к аппаратному обеспечению}
\begin{itemize}
    \item Процессор: x86\_64, 2+ ядра
    \item ОЗУ: 4+ ГБ
    \item Видеокарта с поддержкой OpenGL 3.3+
\end{itemize}

\subsection{Требования к программному обеспечению}
\begin{itemize}
    \item ОС: Linux (Ubuntu 22.04/24.04), Windows 10/11
    \item Python 3.13+, uv
    \item CMake 3.15+ (для сборки C++ модуля)
\end{itemize}

\section{Вызов и загрузка}

\subsection{Установка}
\begin{verbatim}
git clone https://github.com/DIK2022/motion_visualizer
cd motion_visualizer
uv sync
mkdir build && cd build
cmake .. && make
cd ..
\end{verbatim}

\subsection{Запуск}
\begin{verbatim}
uv run python -m motion_visualizer.main
\end{verbatim}

\end{document}